%! AP Statistics — Unit 5, Lesson 5.7 — Homework
\documentclass[10pt]{article}
\usepackage{apstats-article}
\usepackage{apstats-boxes}
\usepackage{grading}
\gradingassignment{apstats-u05-l07-hw}
\gradingcourse{advanced_statistics}

\begin{document}

\pageheader{Unit 5, Lesson 5.7}{Homework}
\namedateperiod

\vspace{0.15in}

\begin{enumerate}

  \item The daily number of steps for adults is right-skewed with $\mu = 7{,}500$ steps
        and $\sigma = 2{,}000$ steps. A researcher surveys $n = 64$ adults at random.
        \begin{enumerate}[label=\alph*.]
          \item Check all three conditions for the normal model for $\bar{x}$.
                \gradedhide{\writeline\\\gradedhere{\gradedblank[id=hw1a, type=rubric, points=3, display={Random sample; $64 \le 10\%$ of all adults; $n = 64 \ge 30$, so the CLT applies despite the skew. Conditions met.}]{\linewidth+6pt}{\gradedlinesheight{2}}}\writelines{2}}
          \item Find $\mu_{\bar{x}}$ and $\sigma_{\bar{x}}$.\gradedbelow[r]{0.5in-0.33in}{$\mu_{\bar{x}} = $~\gradedblank[id=hw1b1, type=short, points=1, answer=7500, check=numeric, tolerance=0]{0.9in}{0.28in}\quad $\sigma_{\bar{x}} = $~\gradedblank[id=hw1b2, type=short, points=1, answer=250, check=numeric, tolerance=0]{0.9in}{0.28in}}
                \vspace{0.5in}
          \item Find $P(\bar{x} < 7{,}000)$.\gradedbelow[r]{0.9in-0.33in}{\gradedblank[id=hw1c, type=short, points=2, answer=0.0228, check=numeric, tolerance=0.002]{0.9in}{0.28in}}
                \vspace{0.9in}
          \item Find $c$ such that $P(\bar{x} > c) = 0.10$. Interpret in context.\gradedbelow[r]{0.9in-0.45in}{$c = $~\gradedblank[id=hw1d1, type=short, points=2, answer=7821, check=numeric, tolerance=2]{0.9in}{0.28in}\quad \gradedblank[id=hw1d2, type=rubric, points=1, display={In 10\% of samples of 64 adults, the mean daily step count exceeds 7,821.}]{3.2in}{0.4in}}
                \vspace{0.9in}
        \end{enumerate}

  \item Exam scores at a large university are normally distributed with $\mu = 74$ and
        $\sigma = 12$. A professor randomly selects $n = 9$ exams to grade first.
        \begin{enumerate}[label=\alph*.]
          \item Are conditions met? Identify each condition.
                \gradedhide{\writeline\\\gradedhere{\gradedblank[id=hw2a, type=rubric, points=3, display={Random sample; $9 \le 10\%$ of all exams; population is normal, so OK for $n = 9$. Met.}]{\linewidth+6pt}{\gradedlinesheight{1}}}\writelines{1}}
          \item Find $P(70 \le \bar{x} \le 78)$.\gradedbelow[r]{0.9in-0.33in}{\gradedblank[id=hw2b, type=short, points=2, answer=0.6826, check=numeric, tolerance=0.003]{0.9in}{0.28in}}
                \vspace{0.9in}
        \end{enumerate}

  \item For a population with $\mu = 200$ and $\sigma = 40$, a researcher takes samples
        of size $n = 25$, $n = 100$, and $n = 400$.
        \begin{enumerate}[label=\alph*.]
          \item Find $\sigma_{\bar{x}}$ for each $n$.\gradedbelow[r]{0.5in-0.33in}{$n = 25$:~\gradedblank[id=hw3a1, type=short, points=1, answer=8, check=numeric, tolerance=0, display={\boldmath$8$}]{0.8in}{0.28in}\quad $n = 100$:~\gradedblank[id=hw3a2, type=short, points=1, answer=4, check=numeric, tolerance=0, display={\boldmath$4$}]{0.8in}{0.28in}\quad $n = 400$:~\gradedblank[id=hw3a3, type=short, points=1, answer=2, check=numeric, tolerance=0, display={\boldmath$2$}]{0.8in}{0.28in}}
                \vspace{0.5in}
          \item For $n = 100$: find $P(195 < \bar{x} < 205)$.\gradedbelow[r]{0.9in-0.33in}{\gradedblank[id=hw3b, type=short, points=2, answer=0.7888, check=numeric, tolerance=0.003]{0.9in}{0.28in}}
                \vspace{0.9in}
        \end{enumerate}

\end{enumerate}

\begin{extensionbox}
\textbf{Extension (optional).}
For $\mu = 200$, $\sigma = 40$, find the smallest $n$ such that $P(|\bar{x} - 200| < 5) \ge 0.95$.
\end{extensionbox}

\begin{reflectionbox}
\textbf{Preview:} Lesson 5.8 extends today's ideas to the difference of two sample means,
$\bar{x}_1 - \bar{x}_2$ --- the foundation for two-sample $t$ procedures in Unit 7.
\end{reflectionbox}

\end{document}
